\documentclass[twocolumn]{jlreq}
\usepackage{luatexja}
\usepackage{luatexja-fontspec}
\usepackage[papersize={148truemm, 100truemm},margin=4pt]{geometry}
\usepackage{fancybox}
\usepackage[style=ieee]{biblatex}
\addbibresource{ref.bib}
\usepackage{xcolor,graphicx,url}
\usepackage{listings,framed}
\usepackage{amsmath,amssymb,mathtools,fontawesome5}
\setsansjfont[
    UprightFont = KouzanMouhituFontOTF,
]{KouzanMouhituFontOTF}
\setsansfont[
    UprightFont = KouzanMouhituFontOTF,
]{KouzanMouhituFontOTF}
\renewcommand\kanjifamilydefault{\gtdefault}
\renewcommand\familydefault{\sfdefault}
\lstset{
    language=C,
    breaklines=true,
    breakindent=10pt,
    basicstyle=\ttfamily\scriptsize,
    frame=TB,
    keywordstyle={\color{red!70!black}\bfseries\ttfamily},
    commentstyle={\color{green!50!black}\ttfamily},
    stringstyle={\color{blue}\ttfamily},
    morekeywords={pipe, fork, pid_t, printf, sprintf, write, read},
    numbers=left,
    numberstyle={\ttfamily\scriptsize},
    numbersep = {1pt}
}

\fancyput(-20truemm,-60truemm){
	\color{black!10}{
		\rotatebox{50}{
			\scalebox{10}{丙午}     
		}
	}
}
\begin{document}
\twocolumn[
    % \vspace{-1cm}
    \begin{center}
        {\Huge 謹賀新年}
    \end{center}
]
{\small
    謹んで新年のご挨拶を申し上げます\par
    旧年中は大変お世話になりありがとうございました\par
    本年も変わらぬご厚誼のほどお願い申し上げます
    \begin{flushright}
        令和7年元旦
    \end{flushright}
}
\vspace{-.3cm}
\begin{framed}
    \vspace{2.5cm}
\end{framed}
% \vspace{-.3cm}
\noindent\scriptsize{\url{https://github.com/MIZOGUCHIKoki/LaTeX/tree/master/new_years_card/2026/main.tex}\\
    \hfill Compiled on Dec. 31, 2025 by Lua\LaTeX}
\vspace{.5cm}
\parbox{.9\linewidth}{
    \setlength{\baselineskip}{5pt}
    {\fontsize{5px}{0px}\selectfont\noindent\input{jusyo.txt}}
    \vspace{.3cm}
    \begin{center}
        {\LARGE 溝\ 口\ \  洸\ 熙}
    \end{center}
    \vspace{.3cm}
}
\newpage
\rmfamily
2026年はハッピー数の年です．\par
ハッピー数とは，各桁の数字を2乗して足し合わせ，その結果に対して同じ操作を繰り返したとき，最終的に1に到達するような数のことを言います．\par
例えば，2026の場合，次のようになります．
\begin{align*}
    \underbrace{2^2 + 0^2 + 2^2 + 6^2}_{2026の各桁の二乗の和} & = 4 + 0 + 4 + 36 = 44 \\
    4^2 + 4^2                                         & = 16 + 16 = 32        \\
    3^2 + 2^2                                         & = 9 + 4 = 13          \\
    1^2 + 3^2                                         & = 1 + 9 = 10          \\
    1^2 + 0^2                                         & = 1 + 0 = 1
\end{align*}
したがって，2026はハッピー数です．\par
皆さんの今年一年が幸せな（\faHandPeace[regular]~~ハッピーな~~\faHandPeace[regular]）年でありますように．

ちなみに，ハッピー数は\(n\)次元に拡張することもできる．\citeauthor{1050564288795241088}による，\(n\)次元ハッピー関数
\begin{align*}
    \mathrm{Happy}(n,x):\mathbb{N}\times \mathbb{N}\cup\{0\}\rightarrow \mathbb{N}\cup\{0\}
\end{align*}
の定義は，文献\cite{1050564288795241088}に詳しい．

\noindent\parbox{\linewidth}{\tiny\cite{1050564288795241088}\fullcite{1050564288795241088}}
\end{document}